\documentclass[11pt]{article}
\usepackage[margin=1in]{geometry}
\usepackage{amsmath,amssymb,amsthm,mathtools,booktabs,microtype}
\usepackage[hidelinks]{hyperref}
\newtheorem{theorem}{Theorem}[section]
\newtheorem{proposition}[theorem]{Proposition}
\newtheorem{corollary}[theorem]{Corollary}
\newtheorem{remark}[theorem]{Remark}
\newcommand{\F}{\mathbb F_2}
\title{Exact Center-Contour Charging on a Four-Dimensional Cubical Block\\
\large The global cost hidden by a frustrated microscopic conditioning}
\author{}
\date{20 July 2026}
\begin{document}
\maketitle

\begin{abstract}
A center-valued boundary conditioning can make one plaquette adjacent to a
selected link frustrated while the other five are aligned.  Treating that
boundary as free hides its global curvature.  We reconstruct the complete
center field on the open side-two four-dimensional cubical block.  The
plaquette complex over $\F$ has ranks $136$ and $80$ on the two adjacent
coboundaries and is exact at all $216$ plaquettes.  An exhaustive syndrome
certificate proves that the one-frustrated/five-aligned incident pattern
requires at least five negative plaquettes globally; a single negative link
attains five.  We then derive a volume-uniform rooted contour bound for the
center subsector.  This is an exact dependent-sector theorem and a calibration
of boundary charging, not control of continuous nonabelian fluctuations.
\end{abstract}

\tableofcontents

\section{Center cochains and Bianchi exactness}

Let $C^1,C^2,C^3$ be the link, plaquette, and three-cell cochains of the open
side-two four-cube, with coefficients in $\F$.  A center-valued link field is
$x\in C^1$; its negative plaquette field is
\[
 y=d_1x\in C^2.
\]
The cubical Bianchi identity is $d_2y=0$.

The block contains $216$ links, $216$ plaquettes, and $96$ three-cells.  Exact
row reduction gives
\begin{equation}
 \operatorname{rank}_{\F}d_1=136,
 \qquad
 \operatorname{rank}_{\F}d_2=80,
 \qquad d_2d_1=0.
 \label{eq:ranks}
\end{equation}

\begin{proposition}[Exactness at plaquettes]
Every Bianchi-closed center plaquette field on this block is an exact link
coboundary:
\[
 \ker d_2=\operatorname{im}d_1.
\]
\end{proposition}
\begin{proof}
The inclusion follows from $d_2d_1=0$.  By \eqref{eq:ranks},
\[
 \dim\ker d_2=216-80=136=\dim\operatorname{im}d_1.
\]
The inclusion is therefore equality.
\end{proof}

\section{Minimum charge of the frustrated incident pattern}

Fix the degree-six link selected in Stage 10.  Among its six incident
plaquettes require one designated plaquette $p_*$ to be negative and the other
five to be positive.  The previously constructed assignment has one negative
center link and full negative-plaquette support
\[
 \{7,17,21,22,23\},
\]
of weight five, with $p_*=21$.

\begin{theorem}[Exact minimum five]
Every center plaquette field satisfying the incident pattern and the Bianchi
identity has Hamming weight at least five.  The bound is attained.
\end{theorem}
\begin{proof}
Represent each of the $216$ columns of $d_2$ as an exact $96$-bit integer
syndrome.  Since $p_*$ is present, all additional plaquettes must have column
XOR equal to the syndrome of $p_*$.  The other five incident plaquettes are
forbidden.  Exact meet-in-the-middle enumeration checks the empty set, every
allowed singleton, every allowed pair, and every allowed triple.  None has the
required XOR.  Thus no Bianchi-closed completion has total weight at most
four.  The displayed five-plaquette coboundary verifies attainment.
\end{proof}

This is a finite exhaustive certificate over integers, rather than a
floating-point optimization.  Independently, no nonzero exact center cocycle
has size one or two: no $d_2$ column is zero and no two columns have the same
syndrome.  A corner link has a three-plaquette coboundary, so the global
nonempty minimum is exactly three.

\section{Center-contour summability}

Join two plaquettes when they are faces of a common three-cell.  This is the
component graph forced by the Bianchi parity constraint: even opposite faces
of one three-cell must remain in the same component.  The exact maximum degree
of this adjacency graph on the block is
\[
 \Delta_{\mathrm B}=20.
\]
Let $q\in(0,1)$ be the relative weight of a negative center plaquette.  In the
center-restricted heat-kernel measure, a curvature support $X$ has relative
weight $q^{|X|}$.  Its connected components are themselves Bianchi cocycles;
on the open block they are exact.

The number of connected plaquette sets of size $n$ containing a fixed root is
bounded by $\Delta_{\mathrm B}^{2(n-1)}$, using a canonical depth-first traversal of a
spanning tree.  Since nonempty center contours have $n\geq3$, we obtain:

\begin{theorem}[Rooted center-contour bound]
If $e^\alpha\Delta_{\mathrm B}^2q<1$, then
\begin{equation}
 \sup_p\sum_{X\ni p}e^{\alpha|X|}|z(X)|
 \leq
 \frac{e^{3\alpha}q^3\Delta_{\mathrm B}^4}{1-e^\alpha\Delta_{\mathrm B}^2q}.
 \label{eq:center-bound}
\end{equation}
The particular Stage 10 incident pattern pays at least $q^5$.
\end{theorem}
\begin{proof}
Discarding the exact-cocycle restriction only enlarges the sum.  Insert the
connected-set count and sum the geometric series from $n=3$.  The exact
minimum-five theorem supplies the final statement.
\end{proof}

\section{Calibration and interpretation}

For $\alpha=0.1$ and $\Delta_{\mathrm B}=20$, the sufficient threshold is
\[
 q<\frac{e^{-0.1}}{20^2}=0.00226209354509.
\]
At the declared calibration $q=10^{-4}$, the branching ratio is
$0.0442068$ and the right side of \eqref{eq:center-bound} is
$2.25965\times10^{-7}$.  The least suppressed realization of the selected
frustration pays $q^5=10^{-20}$.

At small heat time the endpoint heat-kernel ratio has logarithmic rate
$\log q\sim-\pi^2/(2t)$.  For $t=3\times10^{-6}$ this rate proxy is
$-1.644934\times10^6$; prefactors are irrelevant to that asymptotic statement.

\section{Frontier audit}

The theorem demonstrates exactly what arbitrary conditioning erased: a local
frustrated staple belongs to a larger center-curvature contour.  The proof
does not imply that every continuous $\mathrm{SU}(2)$ collar failure pays five
near-antipodal plaquettes.  Continuous fields can spread curvature, and
noncentral orbit/topology labels require separate estimates.  The strongest
claim is C2 for the center subsector.

\end{document}
